\begin{summarybox}
	\textbf{\textcolor{CSummary}{一句话核心}}

	合取全真才真，析取一真即真，否定真值相反。

	\medskip
	\textbf{\textcolor{CSummary}{使用条件}}
	\begin{itemize}[leftmargin=*, itemsep=0.5em, topsep=0.4em]
		\item \textbf{条件 1（命题可判断）：} $p,q$ 必须是能判断真假的命题。
		\item \textbf{条件 2（真值二分）：} 真值仅取 $1$（真）或 $0$（假）。
	\end{itemize}

	\medskip
	\textbf{\textcolor{CSummary}{关键提醒}}
	\begin{itemize}[leftmargin=*, itemsep=0.5em, topsep=0.4em]
		\item \textbf{易错点（区分合取析取）：} 注意“且”与“或”的不同条件。
		\item \textbf{检查项（枚举完整）：} 列真值表务必包含全部 $2^n$ 种赋值。
	\end{itemize}
\end{summarybox}
